% !TEX program = pdflatex
% Example: Response to reviewers letter
% Compile with: pdflatex main.tex

\documentclass[11pt, a4paper]{article}
\usepackage[margin=2.5cm]{geometry}
\usepackage{xcolor}
\usepackage{enumitem}
\usepackage{hyperref}

\newcommand{\reviewer}[1]{\section*{Reviewer #1}}
\newcommand{\comment}[1]{\noindent\textbf{Comment:} #1\par}
\newcommand{\response}[1]{\noindent\textbf{Response:} #1\par\vspace{1em}}

\begin{document}

\title{Response to Reviewers\\[0.5em]
\large Manuscript ID: JSP-2024-0123}
\author{Author Name}
\date{\today}
\maketitle

We thank the reviewers for their constructive feedback.
Below, we address each comment point by point.
Changes in the revised manuscript are highlighted in blue.

\reviewer{1}

\comment{The methodology section lacks detail on the
preprocessing steps. Please clarify how the signals
were filtered and normalized.}

\response{We have expanded Section 3.1 to include a
detailed description of the preprocessing pipeline.
Specifically, we now describe: (1) bandpass filtering
at 0.5--50 Hz using a 4th-order Butterworth filter,
(2) z-score normalization per channel, and (3) artifact
rejection using threshold-based amplitude criteria.
The revised text is on page 5, lines 120--140.}

\comment{Figure 2 is difficult to read. Please improve
the resolution and font sizes.}

\response{We have regenerated Figure 2 at 300 DPI with
larger font sizes (12pt instead of 8pt) for better
readability. We also added axis labels that were
missing in the original submission.}

\reviewer{2}

\comment{The comparison with state-of-the-art methods
is incomplete. Please add comparison with the method
proposed in [Smith et al., 2023].}

\response{We thank the reviewer for this suggestion.
We have added a comparison with Smith et al. (2023)
in Table 2. Our method achieves 93.7\% accuracy
compared to 91.2\% for Smith et al., while using
40\% fewer parameters. The discussion of this
comparison is in Section 4.3, page 9.}

\comment{The statistical significance of the results
is not reported. Please include p-values or confidence
intervals.}

\response{We have added 95\% confidence intervals to
all results in Table 2, computed using bootstrap
resampling with 1000 iterations. We also added
paired t-test p-values comparing our method with
each baseline. All comparisons are statistically
significant (p < 0.01).}

\reviewer{3}

\comment{The paper would benefit from a discussion
of limitations and future work.}

\response{We have added a new Section 5.4 ``Limitations
and Future Work'' that discusses: (1) the computational
cost of the multi-scale architecture, (2) the need for
larger datasets for generalization, and (3) potential
extensions to real-time applications.}

\vspace{1em}
\noindent\textbf{Summary of changes:}
\begin{itemize}
\item Expanded preprocessing description (Section 3.1).
\item Improved Figure 2 resolution and labels.
\item Added comparison with Smith et al. (2023) in Table 2.
\item Added confidence intervals and p-values to Table 2.
\item New Section 5.4 on limitations and future work.
\end{itemize}

\end{document}
